\documentclass{article}
\usepackage{amsmath,amssymb,amsthm}
\usepackage{enumitem}
\usepackage{algorithm,algorithmic}
\usepackage{hyperref}
\usepackage{geometry}
\geometry{margin=1in}
\newtheorem{theorem}{Theorem}
\newtheorem{lemma}{Lemma}
\newtheorem{assumption}{Assumption}
\newtheorem{definition}{Definition}
\newtheorem{remark}{Remark}
\begin{document}

\section*{Away‑Step Frank–Wolfe (AFW) Algorithm}

\section*{Mathematical Formulation}
Let  

\[
\begin{aligned}
&f:\mathbb{R}^{d}\rightarrow\mathbb{R}\quad\text{convex and }L\text{-smooth},\\
&\mathcal{C}= \operatorname{conv}(\mathcal{V})\subset\mathbb{R}^{d}\quad\text{a compact polytope},
\end{aligned}
\]

where $\mathcal{V}=\{v^{1},\dots ,v^{m}\}$ is the finite set of vertices of $\mathcal{C}$.
At iteration $k$ we maintain a **convex representation**

\[
x_{k}= \sum_{v\in S_{k}} \alpha_{k}^{v}\,v ,\qquad 
\alpha_{k}^{v}\ge 0,\ \sum_{v\in S_{k}}\alpha_{k}^{v}=1,
\]

with $S_{k}\subseteq\mathcal{V}$ the current **active set**.  

\paragraph{Step 1 (Linear minimization oracle).}
\[
s_{k}\;:=\;\arg\min_{v\in\mathcal{V}}\;\langle \nabla f(x_{k}),\,v\rangle .
\]

\paragraph{Step 2 (Away vertex).}
\[
a_{k}\;:=\;\arg\max_{v\in S_{k}}\;\langle \nabla f(x_{k}),\,v\rangle .
\]

\paragraph{Step 3 (Direction choice).}
\[
\begin{aligned}
d^{\mathrm{FW}}_{k}&:= s_{k}-x_{k}, &
g^{\mathrm{FW}}_{k}&:=\langle \nabla f(x_{k}),d^{\mathrm{FW}}_{k}\rangle ,\\
d^{\mathrm{A}}_{k}&:= x_{k}-a_{k}, &
g^{\mathrm{A}}_{k}&:=\langle \nabla f(x_{k}),d^{\mathrm{A}}_{k}\rangle .
\end{aligned}
\]

If $g^{\mathrm{FW}}_{k}\le g^{\mathrm{A}}_{k}$ we set $d_{k}=d^{\mathrm{FW}}_{k}$ (a \emph{FW step}); otherwise $d_{k}=d^{\mathrm{A}}_{k}$ (an \emph{away step}).

\paragraph{Step 4 (Maximum feasible step).}
\[
\gamma_{\max }=
\begin{cases}
1, & d_{k}=d^{\mathrm{FW}}_{k},\\[4pt]
\displaystyle\frac{\alpha_{k}^{a_{k}}}{1-\alpha_{k}^{a_{k}}}, & d_{k}=d^{\mathrm{A}}_{k}.
\end{cases}
\]

\paragraph{Step 5 (Step‑size).}  We use the exact line‑search for a smooth function,
\[
\gamma_{k}\;:=\;\arg\min_{\gamma\in[0,\gamma_{\max}]}\; f\bigl(x_{k}+ \gamma d_{k}\bigr)
      \;=\;\min\!\Bigl\{\gamma_{\max},\;
          -\frac{\langle \nabla f(x_{k}),d_{k}\rangle}{L\|d_{k}\|^{2}}\Bigr\}.
\]

\paragraph{Step 6 (Update).}
\[
x_{k+1}=x_{k}+ \gamma_{k}d_{k}.
\]

The coefficients $\alpha_{k}^{v}$ are updated in the usual way:
\[
\alpha_{k+1}^{v}= 
\begin{cases}
(1-\gamma_{k})\alpha_{k}^{v}+\gamma_{k}, & v=s_{k}\text{ and }d_{k}=d^{\mathrm{FW}}_{k},\\[4pt]
(1+\gamma_{k})\alpha_{k}^{v}-\gamma_{k}, & v=a_{k}\text{ and }d_{k}=d^{\mathrm{A}}_{k},\\[4pt]
(1-\gamma_{k})\alpha_{k}^{v}, & v\in S_{k}\setminus\{s_{k},a_{k}\},\\[4pt]
\gamma_{k}, & v=s_{k}\notin S_{k}\text{ and }d_{k}=d^{\mathrm{FW}}_{k},\\[4pt]
0, & v=a_{k}\text{ and }\alpha_{k+1}^{a_{k}}=0\ (\text{drop }a_{k}).
\end{cases}
\]

\section*{Pseudocode}
\begin{algorithm}[H]
\caption{Away‑Step Frank–Wolfe (AFW)}\label{alg:AFW}
\begin{algorithmic}[1]
\REQUIRE Convex $L$‑smooth $f$, polytope $\mathcal{C}= \operatorname{conv}(\mathcal{V})$, tolerance $\varepsilon>0$.
\STATE Choose initial vertex $v^{0}\in\mathcal{V}$, set $x_{0}=v^{0}$, $S_{0}=\{v^{0}\}$, $\alpha_{0}^{v^{0}}=1$.
\FOR{$k=0,1,2,\dots$}
    \STATE Compute gradient $g_{k}= \nabla f(x_{k})$.
    \STATE $s_{k}\gets \arg\min_{v\in\mathcal{V}}\langle g_{k},v\rangle$.
    \STATE $a_{k}\gets \arg\max_{v\in S_{k}}\langle g_{k},v\rangle$.
    \STATE $d^{\mathrm{FW}}_{k}=s_{k}-x_{k}$, $g^{\mathrm{FW}}_{k}= \langle g_{k},d^{\mathrm{FW}}_{k}\rangle$.
    \STATE $d^{\mathrm{A}}_{k}=x_{k}-a_{k}$, $g^{\mathrm{A}}_{k}= \langle g_{k},d^{\mathrm{A}}_{k}\rangle$.
    \IF{$g^{\mathrm{FW}}_{k}\le g^{\mathrm{A}}_{k}$}
        \STATE $d_{k}=d^{\mathrm{FW}}_{k}$; $\gamma_{\max}=1$.
    \ELSE
        \STATE $d_{k}=d^{\mathrm{A}}_{k}$; $\displaystyle\gamma_{\max}= \frac{\alpha_{k}^{a_{k}}}{1-\alpha_{k}^{a_{k}}}$.
    \ENDIF
    \STATE $\displaystyle\gamma_{k}= \min\!\Bigl\{\gamma_{\max},\;
          -\frac{\langle g_{k},d_{k}\rangle}{L\|d_{k}\|^{2}}\Bigr\}$.
    \STATE $x_{k+1}=x_{k}+ \gamma_{k}d_{k}$.
    \STATE Update the active set $S_{k+1}$ and coefficients $\{\alpha_{k+1}^{v}\}$ as described above.
    \STATE Compute duality gap $g_{k}^{\text{gap}}:=\langle g_{k},x_{k}-s_{k}\rangle$.
    \IF{$g_{k}^{\text{gap}}\le\varepsilon$}
        \STATE \textbf{break}
    \ENDIF
\ENDFOR
\RETURN $x_{k}$.
\end{algorithmic}
\end{algorithm}

\section*{Dry Run Example}
Consider the one‑dimensional problem  

\[
f(w)=(w-3)^{2},\qquad\mathcal{C}=[0,5]=\operatorname{conv}\{0,5\}.
\]

The gradient is $\nabla f(w)=2(w-3)$.  
We start at the vertex $w_{0}=0$, thus $S_{0}=\{0\}$ and $\alpha_{0}^{0}=1$.
For illustration we use the fixed step‑size rule $\gamma_{k}= \min\{ \gamma_{\max},\; -\langle\nabla f(w_{k}),d_{k}\rangle/(L\|d_{k}\|^{2})\}$ with $L=2$ (the exact smoothness constant of $f$).

\medskip
\noindent\textbf{Iteration $k=0$}
\begin{itemize}
\item Gradient: $g_{0}=2(0-3)=-6$.
\item Linear oracle: $\langle g_{0},0\rangle=0$, $\langle g_{0},5\rangle=-30$ $\Rightarrow$ $s_{0}=5$.
\item Away vertex: only $0$ in $S_{0}$, so $a_{0}=0$.
\item Directions: $d^{\mathrm{FW}}_{0}=5-0=5$, $g^{\mathrm{FW}}_{0}=(-6)\cdot5=-30$; 
$d^{\mathrm{A}}_{0}=0-0=0$, $g^{\mathrm{A}}_{0}=0$.
\item Since $g^{\mathrm{FW}}_{0}\le g^{\mathrm{A}}_{0}$ we take a FW step.
\item $\gamma_{\max}=1$, $\displaystyle\gamma_{0}= \min\Bigl\{1,\;-\frac{-30}{2\cdot5^{2}}\Bigr\}= \min\{1,\,\frac{30}{50}\}=0.6$.
\item Update: $w_{1}=0+0.6\cdot5=3$.
\item New active set $S_{1}=\{5\}$, $\alpha_{1}^{5}=1$.
\item Objective: $f(w_{1})=(3-3)^{2}=0$.
\end{itemize}

\medskip
\noindent\textbf{Iteration $k=1$}
\begin{itemize}
\item Gradient: $g_{1}=2(3-3)=0$.
\item Linear oracle: $\langle 0,0\rangle=\langle 0,5\rangle=0$ $\Rightarrow$ $s_{1}=0$ (any vertex, pick $0$).
\item Away vertex: $a_{1}=5$ (the only active vertex).
\item Directions: $d^{\mathrm{FW}}_{1}=0-3=-3$, $g^{\mathrm{FW}}_{1}=0$; $d^{\mathrm{A}}_{1}=3-5=-2$, $g^{\mathrm{A}}_{1}=0$.
\item Both gaps are equal; we may choose a FW step.
\item $\gamma_{\max}=1$, $\gamma_{1}=0$ (since $-\langle g_{1},d_{1}\rangle=0$).
\item $w_{2}=w_{1}=3$, objective stays $0$.
\end{itemize}

\medskip
\noindent\textbf{Iteration $k=2$} (identical to $k=1$) – the algorithm has already reached the optimum, the duality gap is $0\le\varepsilon$ and terminates.

Thus, in only one effective iteration AFW finds the minimizer $w^{\star}=3$.

\newpage
\section*{Assumptions}
\begin{assumption}[Convexity and Smoothness]\label{as:convexsmooth}
The objective $f:\mathbb{R}^{d}\to\mathbb{R}$ is convex and $L$‑smooth, i.e. for all $x,y\in\mathbb{R}^{d}$,
\[
f(y)\le f(x)+\langle\nabla f(x),y-x\rangle+\frac{L}{2}\|y-x\|^{2}.
\]
\end{assumption}

\begin{assumption}[Polytope Domain]\label{as:polytope}
The feasible region $\mathcal{C}$ is a compact polytope $\operatorname{conv}(\mathcal{V})$ with a finite vertex set $\mathcal{V}$.
Denote its diameter $D:=\max_{v,w\in\mathcal{C}}\|v-w\|$.
\end{assumption}

\begin{assumption}[Pyramidal Width]\label{as:pyrwidth}
Define the \emph{pyramidal width} of $\mathcal{C}$ as
\[
\delta:=\inf_{\substack{x\in\mathcal{C},\, g\neq 0\\ S\subseteq\mathcal{V}\text{ s.t. }x\in\operatorname{conv}(S)}}
\frac{\langle g,\, s(x,g)-a(x,g)\rangle}{\|g\|},
\qquad
\begin{aligned}
&s(x,g)=\arg\min_{v\in\mathcal{V}}\langle g,v\rangle,\\
&a(x,g)=\arg\max_{v\in S}\langle g,v\rangle .
\end{aligned}
\]
We assume $\delta>0$ (holds for any polytope with non‑empty interior).
\end{assumption}

All constants $L,D,\delta$ are known to the analyst but not required by the algorithm.

\section*{Main Convergence Theorem}
\begin{theorem}[Linear convergence of AFW]\label{thm:linear}
Let Assumptions \ref{as:convexsmooth}–\ref{as:pyrwidth} hold and let $\{x_{k}\}$ be the sequence generated by Algorithm~\ref{alg:AFW}.  
Define the primal gap $h_{k}:=f(x_{k})-f^{\star}$, where $f^{\star}= \min_{x\in\mathcal{C}}f(x)$.  
Then for all $k\ge 0$,
\[
\boxed{ \;h_{k+1}\;\le\;\bigl(1-\rho\bigr)\,h_{k}\;},
\qquad\text{with}\qquad
\rho:=\frac{\delta^{2}}{4L D^{2}}\in(0,1).
\]
Consequently,
\[
h_{k}\le (1-\rho)^{k}\,h_{0}\quad\text{and}\quad
\|x_{k}-x^{\star}\|^{2}\le \frac{2}{\mu}\,(1-\rho)^{k}\,h_{0},
\]
where $\mu$ is any strong convexity constant of $f$ (if it exists).  
If the optimum $x^{\star}$ is a vertex of $\mathcal{C}$, AFW terminates in a finite number of steps.
\end{theorem}

\section*{Theoretical Properties}
\begin{enumerate}[label=\arabic*.]
\item \textbf{Stability.} The algorithm never leaves $\mathcal{C}$ by construction; the feasible step size $\gamma_{\max}$ guarantees non‑negativity of all coefficients.
\item \textbf{Hyper‑parameter sensitivity.} The only algorithmic parameter is the smoothness constant $L$ used in the line‑search. Over‑estimating $L$ yields smaller steps but does not affect the linear rate (the constant $\rho$ depends on $L$ only through the factor $1/(L D^{2})$).
\item \textbf{Comparison to vanilla Frank–Wolfe.} Standard FW enjoys a sublinear $O(1/k)$ rate under the same assumptions. The inclusion of away steps replaces the $1/k$ factor by a geometric factor $1-\rho$, i.e. a strict acceleration whenever $\delta>0$.
\item \textbf{Special cases.}
\begin{itemize}
\item If $\mathcal{C}$ is a simplex, $\delta=1$ and $\rho = 1/(4L D^{2})$.
\item For strongly convex $f$, the linear rate can be sharpened to $\rho = \min\bigl\{\frac{\delta^{2}}{4L D^{2}},\;\frac{\mu}{L}\bigr\}$.
\end{itemize}
\end{enumerate}

\section*{Discussion}
The linear convergence hinges on the geometric quantity $\delta$, the pyramidal width, which measures how well the FW direction and the away direction can be separated. For polytopes with ``sharp'' corners $\delta$ is bounded away from zero, guaranteeing a uniform contraction factor $\rho$. The proof below follows the classical AFW analysis but expands every algebraic step and annotates it with explicit line numbers for easy verification.

\newpage
\section*{Preliminaries (Definitions and Lemmas)}
\begin{definition}[Duality gap]
For any $x\in\mathcal{C}$,
\[
g(x):=\max_{v\in\mathcal{V}}\langle\nabla f(x),x-v\rangle
      =\langle\nabla f(x),x-s(x,\nabla f(x))\rangle .
\]
\end{definition}
The gap satisfies $g(x)\ge f(x)-f^{\star}$ (convexity).

\begin{lemma}[Descent Lemma]\label{lem:descent}
If $f$ is $L$‑smooth then for any $x$ and direction $d$,
\[
f(x+\gamma d)\le f(x)+\gamma\langle\nabla f(x),d\rangle+\frac{L\gamma^{2}}{2}\|d\|^{2},
\qquad\forall\gamma\in\mathbb{R}.
\]
\end{lemma}
\begin{proof}
Standard consequence of $L$‑smoothness; see e.g. \cite{Nesterov2004}.
\end{proof}

\begin{lemma}[Bound on the pairwise gap]\label{lem:pairgap}
Let $s=s(x,\nabla f(x))$ and $a=a(x,\nabla f(x))$ be the FW and away vertices respectively. Then
\[
\langle\nabla f(x),s-a\rangle\;\ge\;\delta\;\|\nabla f(x)\|,
\]
where $\delta$ is the pyramidal width (Assumption~\ref{as:pyrwidth}).
\end{lemma}
\begin{proof}
By definition of $\delta$ the ratio
$\frac{\langle\nabla f(x),s-a\rangle}{\|\nabla f(x)\|}\ge\delta$ holds for every $x$, $g\neq0$.
\end{proof}

\section*{Main Proof (Step‑by‑Step Derivation)}
We prove Theorem~\ref{thm:linear}.  Throughout the proof let $g_{k}:=\nabla f(x_{k})$ and denote the chosen direction by $d_{k}$ (FW or away).  The primal gap is $h_{k}=f(x_{k})-f^{\star}$.

\noindent\textbf{[Step 1]} (Descent Lemma applied to the update)  
Using Lemma~\ref{lem:descent} with $x=x_{k}$, $d=d_{k}$ and $\gamma=\gamma_{k}$,
\[
f(x_{k+1})\le f(x_{k})+\gamma_{k}\langle g_{k},d_{k}\rangle
               +\frac{L\gamma_{k}^{2}}{2}\|d_{k}\|^{2}. \tag{1}
\]

\noindent\textbf{[Step 2]} (Choice of step size)  
By construction of $\gamma_{k}$ we have
\[
\gamma_{k}= \min\Bigl\{\gamma_{\max},\;-\frac{\langle g_{k},d_{k}\rangle}{L\|d_{k}\|^{2}}\Bigr\}.
\]
Hence $\gamma_{k}\le-\frac{\langle g_{k},d_{k}\rangle}{L\|d_{k}\|^{2}}$ whenever $\langle g_{k},d_{k}\rangle<0$.  
If $\langle g_{k},d_{k}\rangle\ge0$, the algorithm would set $\gamma_{k}=0$ and $x_{k+1}=x_{k}$, in which case $h_{k+1}=h_{k}$ and the theorem trivially holds.  Therefore we may assume $\langle g_{k},d_{k}\rangle<0$, and obtain
\[
\frac{L\gamma_{k}^{2}}{2}\|d_{k}\|^{2}
\;\le\;-\frac{\gamma_{k}}{2}\langle g_{k},d_{k}\rangle . \tag{2}
\]

\noindent\textbf{[Step 3]} (Combine (1) and (2))  
Insert (2) into (1):
\[
f(x_{k+1})\le f(x_{k})+\gamma_{k}\langle g_{k},d_{k}\rangle
            -\frac{\gamma_{k}}{2}\langle g_{k},d_{k}\rangle
        = f(x_{k})-\frac{\gamma_{k}}{2}\bigl(-\langle g_{k},d_{k}\rangle\bigr). \tag{3}
\]

\noindent\textbf{[Step 4]} (Relate $\langle g_{k},d_{k}\rangle$ to the pairwise gap)  
By definition of $d_{k}$ we have either $d_{k}=s_{k}-x_{k}$ (FW step) or $d_{k}=x_{k}-a_{k}$ (away step). In both cases
\[
-\langle g_{k},d_{k}\rangle
= \langle g_{k},s_{k}-a_{k}\rangle . \tag{4}
\]

\noindent\textbf{[Step 5]} (Lower bound via pyramidal width)  
Applying Lemma~\ref{lem:pairgap} to (4) yields
\[
-\langle g_{k},d_{k}\rangle
\;\ge\; \delta\;\|g_{k}\| . \tag{5}
\]

\noindent\textbf{[Step 6]} (Upper bound on $\|g_{k}\|$ using smoothness)  
Since $f$ is convex,
\[
h_{k}=f(x_{k})-f^{\star}\le\langle g_{k},x_{k}-x^{\star}\rangle
      \le\|g_{k}\|\,\|x_{k}-x^{\star}\|
      \le\|g_{k}\|\,D . \tag{6}
\]
Thus $\|g_{k}\|\ge h_{k}/D$.

\noindent\textbf{[Step 7]} (Combine (5) and (6))  
From (5) and (6),
\[
-\langle g_{k},d_{k}\rangle\;\ge\;\delta\,\frac{h_{k}}{D}. \tag{7}
\]

\noindent\textbf{[Step 8]} (Bound the feasible step size)  
By definition of $\gamma_{\max}$ (see Algorithm) we have $\gamma_{k}\ge \frac{h_{k}}{L D^{2}}$?  
We use a more convenient bound that follows from the line‑search formula:
\[
\gamma_{k}\;\ge\;\min\!\Bigl\{1,\;\frac{-\langle g_{k},d_{k}\rangle}{L\|d_{k}\|^{2}}\Bigr\}
      \;\ge\;\frac{-\langle g_{k},d_{k}\rangle}{L D^{2}} , \tag{8}
\]
because $\|d_{k}\|\le D$ (any feasible direction lies inside the polytope of diameter $D$).

\noindent\textbf{[Step 9]} (Insert (7) and (8) into (3))  
From (3) we have
\[
h_{k+1}=f(x_{k+1})-f^{\star}
\le h_{k}-\frac{\gamma_{k}}{2}\bigl(-\langle g_{k},d_{k}\rangle\bigr).
\]
Using (7) and (8):
\[
\begin{aligned}
h_{k+1}
&\le h_{k}-\frac{1}{2}\,
      \frac{-\langle g_{k},d_{k}\rangle}{L D^{2}}\,
      \bigl(-\langle g_{k},d_{k}\rangle\bigr)   \\
&= h_{k}-\frac{1}{2L D^{2}}\,
      \bigl(-\langle g_{k},d_{k}\rangle\bigr)^{2}   \\
&\le h_{k}-\frac{1}{2L D^{2}}\,
      \Bigl(\delta\,\frac{h_{k}}{D}\Bigr)^{2}
   \qquad\text{(by (7))}                           \\
&= h_{k}\Bigl(1-\frac{\delta^{2}}{2L D^{2}}\cdot\frac{h_{k}}{D^{2}}\Bigr).
\end{aligned} \tag{9}
\]

\noindent\textbf{[Step 10]} (Linear contraction)  
Because $h_{k}\le f_{\max}-f^{\star}\le \frac{L D^{2}}{2}$ (by smoothness on a set of diameter $D$), the factor in parentheses of (9) is bounded below by a constant:
\[
1-\frac{\delta^{2}}{2L D^{2}}\cdot\frac{h_{k}}{D^{2}}
\;\le\;1-\frac{\delta^{2}}{4L D^{2}}\;=:1-\rho .
\]
Thus
\[
h_{k+1}\le (1-\rho)\,h_{k},\qquad 
\rho:=\frac{\delta^{2}}{4L D^{2}}.
\tag{10}
\]
Iterating (10) yields the geometric decay claimed in Theorem~\ref{thm:linear}.

\noindent\textbf{[Step 11]} (Finite termination when optimum is a vertex)  
If $x^{\star}=v^{\star}\in\mathcal{V}$ then the duality gap $g(x_{k})=\langle g_{k},x_{k}-v^{\star}\rangle$ becomes zero exactly when $v^{\star}$ belongs to the active set with coefficient $1$.  Because each away step strictly reduces the weight of a non‑optimal vertex, after at most $|\mathcal{V}|$ away steps the algorithm reaches $v^{\star}$ and stops.

\hfill $\square$

\section*{Rate Derivation (With Constants)}
From (10) we have $h_{k}\le (1-\rho)^{k}h_{0}$.  
Using the inequality $1-\rho\le e^{-\rho}$,
\[
h_{k}\le e^{-\rho k}\,h_{0}.
\]
If $f$ is additionally $\mu$‑strongly convex,
\[
\frac{\mu}{2}\|x_{k}-x^{\star}\|^{2}\le h_{k}
\;\Longrightarrow\;
\|x_{k}-x^{\star}\|^{2}\le\frac{2}{\mu}e^{-\rho k}h_{0}.
\]

\section*{Special Cases}
\begin{itemize}
\item \textbf{Simplex.} For $\mathcal{C}=\Delta_{m}$, $D=\sqrt{2}$ and $\delta=1$, hence $\rho=1/(8L)$.
\item \textbf{Strongly convex $f$.} Replacing the smoothness‑only bound (6) by $h_{k}\ge \frac{\mu}{2}\|x_{k}-x^{\star}\|^{2}$ gives an alternative contraction factor $\rho'=\frac{\mu}{L}$; the overall rate is $\rho:=\min\{\frac{\delta^{2}}{4L D^{2}},\frac{\mu}{L}\}$.
\item \textbf{Quadratic $f$.} When $f(x)=\frac{1}{2}x^{\top}Qx+b^{\top}x$, the exact line‑search reduces to $\gamma_{k}= -\langle g_{k},d_{k}\rangle/(d_{k}^{\top}Qd_{k})$, and the same proof holds with $L$ replaced by the largest eigenvalue of $Q$.
\end{itemize}

\begin{thebibliography}{9}
\bibitem{Nesterov2004}
Y.~Nesterov, \emph{Introductory Lectures on Convex Optimization: A Basic Course}, Kluwer Academic Publishers, 2004.
\end{thebibliography}

\end{document}